\subsubsection{Extension Theorems}\label{subsec:sobolev-extension-theorems}

We now address the question whether a function in $H^k(\Omega)$ can be extended to a function in $H^k(\mathbb{R}^m)$. We first note a trivial extension result.

\begin{Lemma}
\label{lem:zero-boundary-sobolev-zero-extension}
Let $u\in H_0^k(\Omega)$. Then $u$ can be extended to a function in $H^k(\mathbb{R}^m)$ by defining it to be zero outside $\overline\Omega$, i.e., we set $(u,\phi)=\int_{\Omega}\overline{u(\bm{x})}\phi(\bm{x})\,\mathrm{d}\bm{x}$ for every $\phi\in\mathcal{D}(\mathbb{R}^m)$.
\end{Lemma}

\begingroup
\postdisplaypenalty=10000
\renewcommand{\endProof}{\par\nobreak\hfill$\blacksquare$\par\vspace{0.1cm}}
\begin{Proof}
Take a sequence $u_n\in\mathcal{D}(\Omega)$ which converges to $u$ in $H^k(\Omega)$. Then $u_n$ also converges in $H^k(\mathbb{R}^m)$.
\end{Proof}
\endgroup

We next consider the case of a half-space. Let $\mathbb{R}_+^m=\{x_m>0\}$. We have the following result.

\begin{Theorem}
\label{thm:halfspace-sobolev-bounded-extension}
Let $k\geq0$ be an integer. Then there exists a bounded linear mapping $E:H^k(\mathbb{R}_+^m)\to H^k(\mathbb{R}^m)$ with the property that $(Eu)|_{\mathbb{R}_+^m}=u$ for every $u\in H^k(\mathbb{R}_+^m)$. Moreover, for any given $K\in\mathbb{N}$, $E$ can be chosen independently of $k$ for $0\leq k\leq K$.
\end{Theorem}

\begingroup
\postdisplaypenalty=10000
\renewcommand{\endProof}{\par\nobreak\hfill$\blacksquare$\par\vspace{0.1cm}}
\begin{Proof}
Let us first consider how we would extend smooth functions. A continuous function can be extended to a continuous function on the whole space by even continuation, i.e., $u(\bm{x}^{\prime},x_m)=u(\bm{x}^{\prime},-x_m)$ for $x_m<0$. Here and in the following, $\bm{x}^{\prime}$ stands for $(x_1,x_2,\ldots,x_{m-1})$. However, even continuation will generally lead to discontinuous derivatives. If we set $u(\bm{x}^{\prime},x_m)=4u(\bm{x}^{\prime},-x_m/2)-3u(\bm{x}^{\prime},-x_m)$ for $x_m<0$, then both $u$ and $\partial u/\partial x_m$ turn out to be continuous across $x_m=0$ if $u$ is smooth. In general, we make the ansatz
\begin{equation}\label{eq:halfspace-extension-reflected-linear-combination}
u(\bm{x}^{\prime},x_m)=\sum_{j=0}^K\alpha_j u\left(\bm{x}^{\prime},-\dfrac{x_m}{j+1}\right)
\end{equation}
for $x_m<0$, and we want to choose the $\alpha_j$ such that for smooth $u$ the derivatives up to order $K$ match across $x_m=0$. This leads to the equations
\begin{equation}\label{eq:halfspace-extension-vandermonde-matching-system}
\sum_{j=0}^K\left(-\dfrac{1}{j+1}\right)^i\alpha_j=1,\quad i=0,1,\ldots,K.
\end{equation}
The matrix of this system is known as a Vandermonde matrix and is well known to be nonsingular (Problem~\ref{ex:vandermonde-matrix-nonsingularity}). Hence \eqref{eq:halfspace-extension-vandermonde-matching-system} is uniquely solvable.

It is easy to see that if $u\in C^k(\overline{\mathbb{R}_+^m})$ $(k\leq K)$ with bounded support, then $Eu\in C^k(\mathbb{R}^m)$ with bounded support, and $\|Eu\|_{k,2}\leq C\|u\|_{k,2}$ for some constant $C$ that is independent of $u$. If we can show that functions with bounded support in $C^k(\overline{\mathbb{R}_+^m})$ are dense in $H^k(\mathbb{R}_+^m)$, we are done.
\end{Proof}
\endgroup

\begin{Lemma}
\label{lem:halfspace-restricted-test-function-density}
Let $u\in H^k(\mathbb{R}_+^m)$. Then there is a sequence $\phi_n\in\mathcal{D}(\mathbb{R}^m)$ such that the restrictions of $\phi_n$ to $\mathbb{R}_+^m$ converge to $u$.
\end{Lemma}

\begingroup
\postdisplaypenalty=10000
\renewcommand{\endProof}{\par\nobreak\hfill$\blacksquare$\par\vspace{0.1cm}}
\begin{Proof}
Since we already know that $\mathcal{D}(\mathbb{R}^m)$ is dense in $H^k(\mathbb{R}^m)$, it suffices to show that every $u\in H^k(\mathbb{R}_+^m)$ can be approximated by restrictions of functions in $H^k(\mathbb{R}^m)$. Moreover, we know functions of bounded support are dense; hence let us assume $u$ has bounded support. Let $u_n(\bm{x}^{\prime},x_m)=u(\bm{x}^{\prime},x_m+1/n)$ for $\bm{x}\in\mathbb{R}_+^m$. Then $\|u_n\|_2\leq\|u\|_2$, and $u_n\to u$ for $u\in\mathcal{D}(\mathbb{R}_+^m)$. By density, $u_n\to u$ in $L^2$ for every $u\in L^2(\mathbb{R}_+^m)$. Since, moreover, $D^\alpha u_n(\bm{x}^{\prime},x_m)=D^\alpha u(\bm{x}^{\prime},x_m+1/n)$, the same argument shows that $u_n\to u$ in $H^k(\mathbb{R}_+^m)$. We claim that $u_n$ can be extended. First we note that $u_n$ is actually defined for $x_m>-1/n$. Let
\begin{equation}\label{eq:halfspace-density-shifted-zero-extension-distribution}
(v_n,\phi)=\int_{x_m>-1/n}\overline{u_n(\bm{x})}\phi(\bm{x})\,\mathrm{d}\bm{x}
\end{equation}
be the extension of $u_n$ by zero. Then clearly $D^\alpha v_n=0$ for $x_m<-1/n$ and $D^\alpha v_n=D^\alpha u_n$ for $x_m>-1/n$. However, we cannot claim that $D^\alpha v_n$ is the extension of $D^\alpha u_n$ by zero. What we can claim is that
\begin{equation}\label{eq:halfspace-density-zero-extension-boundary-distribution}
(D^\alpha v_n,\phi)=\int_{x_m>-1/n}\overline{D^\alpha u_n(\bm{x})}\phi(\bm{x})\,\mathrm{d}\bm{x}+(g_{n\alpha},\phi),
\end{equation}
where $g_{n\alpha}$ is a distribution with bounded support contained in the plane $x_m=-1/n$. By Lemma~\ref{lem:distribution-local-finite-order}, $g_{n\alpha}$ is of finite order. Let now $\psi_n(x_m)$ be a $C^{\infty}$-function which vanishes for $x_m\leq-1/n$ and equals $1$ for $x_m\geq0$. Then $g_{n\alpha}\psi_n$ vanishes, and hence $w_n=v_n\psi_n$ is in $H^k(\mathbb{R}^m)$. It is clear that the restriction of $w_n$ to $\mathbb{R}_+^m$ is $u_n$. This concludes the proof.
\end{Proof}
\endgroup
